% Laget automatisk av elevark/lag_elevark.py fra sim/vognstot.html. Endringer her blir overskrevet.
\documentclass[11pt]{article}
\usepackage{elevark}
\begin{document}
\elevarkhode{Fysikk 1 · Bevegelsesmengde og støt}{Støt mellom to vogner}{https://fysikklab.no/sim/vognstot.html}{Simuleringen}

To vogner ruller mot hverandre på en vannrett bane uten friksjon. De kan ha magneter som spretter dem fra hverandre, borrelås som får dem til å henge sammen, eller en fjær som skyter dem fra hverandre. Hva er bevart i støtet, og hva er ikke bevart?

\section*{Oppgaver}

Oppgave 1–6 bruker simuleringen. Oppgave 7 bruker egne målinger fra forsøket med vogner.

\begin{oppgaveliste}
\oppgave \fase{gjett}{Gjett}Vogn A (0,50~kg) kjører med 0,60~m/s mot vogn B (0,25~kg), som står i ro. Støtet er elastisk. Fortsetter A framover, stopper den, eller spretter den tilbake?
\svar
\oppgave \fase{test}{Test}Kjør simuleringen og se på grafen for bevegelsesmengde. Hva skjer med $p_{\mathrm{A}}$, $p_{\mathrm{B}}$ og summen under støtet? Bytt så til kinetisk energi. Hvorfor har den samlede kinetiske energien en dupp midt i støtet, selv om støtet er elastisk?
\svar
\oppgave \fase{test}{Regn}Bytt til borrelås med de samme startverdiene. Regn ut felles fart etter støtet og hvor mye kinetisk energi som går tapt. Sjekk med simuleringen.
\svar
\oppgave \fase{forklar}{Forklar}Se på grafen for kraft mellom vognene. Hvordan henger kraften på A sammen med kraften på B? Bruk dette og Newtons 2. lov på formen $ΣF = Δp/Δt$ til å forklare hvorfor den samlede bevegelsesmengden er bevart.
\svar
\oppgave \fase{gjett}{Gjett}Velg eksplosjon med $m_{\mathrm{A}} = 0{,}25\,\mathrm{kg}$ og $m_{\mathrm{B}} = 0{,}75\,\mathrm{kg}$. Hvilken vogn får størst fart? Hvilken får størst bevegelsesmengde? Hvilken får mest kinetisk energi?
\svar
\oppgave \fase{test}{Undersøk}Prøv deg fram: Hvilke masser og farter gjør at vogn A stopper helt etter et elastisk støt? Og hvordan får du begge vognene til å stå i ro etter et støt med borrelås?
\svar
\oppgave \fase{utfordring}{Forsøk}Gjør forsøket med ekte vogner og fyll inn tabellen under «Egne målinger». Er bevegelsesmengden bevart innenfor måleusikkerheten? Hvilke av støtene er tilnærmet elastiske? Hvilke feilkilder kan forklare avvikene?
\svar
\end{oppgaveliste}

\end{document}
